\section{同伦不变性}
\label{sec:spectral-homotopy-invariance}

下述命题说明函数 $\rho_a=\rho(\,\cdot\,;a)$ 下降为
$\widetilde{\operatorname{Ham}}(M,\omega)$ 上的连续函数。

\begin{theorem}[同伦不变性]
\label{thm:21-spectral-homotopy-invariance}
设 $(M,\omega)$ 是任意闭辛流形，并假设 $(M,\omega)$ 满足非退化谱性公理。则对任意
光滑函数 $H\sim K$，即满足 $\phi_H^1=\phi_K^1$ 且 $\phi_H,\phi_K$ 相对于端点道路
同伦的函数，都有
\begin{equation}
 \rho(\underline H;a)=\rho(\underline K;a).
 \label{eq:21-spectral-homotopy-equality}
\end{equation}
\end{theorem}

本节将给出该定理的两个不同证明，每一个都包含弗洛尔理论研究中的若干重要成分。
暂且推迟证明，先陈述其推论。由同伦不变性，可对代表
$h=[\phi_H]$ 的任意 $H$（无论 $h$ 是否非退化）定义
\begin{equation}
 \rho_a(h):=\rho(H;a).
 \label{eq:21-spectral-function-cover}
\end{equation}
这给出良定义函数
\[
 \rho_a:\widetilde{\operatorname{Ham}}(M,\omega)\longrightarrow\mathbb R.
\]

\begin{theorem}
\label{thm:21-spectral-function-continuous-cover}
在由 $P^{\mathrm{ham}}(\operatorname{Symp}(M,\omega),\mathrm{id})$ 诱导的商拓扑下，
$\rho_a$ 是 $\widetilde{\operatorname{Ham}}(M,\omega)$ 上的连续函数。
\end{theorem}

\begin{proof}
回忆商映射
\[
 \pi:P^{\mathrm{ham}}(\operatorname{Symp}(M,\omega),\mathrm{id})
 \longrightarrow\widetilde{\operatorname{Ham}}(M,\omega).
\]
已经证明赋值
\begin{equation}
 \lambda\longmapsto\rho(\lambda;a)
 \label{eq:21-spectral-function-path-space}
\end{equation}
在道路空间上连续。按商拓扑定义，
$\rho_a:\widetilde{\operatorname{Ham}}(M,\omega)\to\mathbb R$ 连续，因为复合
\[
 \rho_a\circ\pi:
 P^{\mathrm{ham}}(\operatorname{Symp}(M,\omega),\mathrm{id})\longrightarrow\mathbb R
\]
正是~\eqref{eq:21-spectral-function-path-space}。
\end{proof}

\subsection{使用非退化谱性的证明}

先考虑满足 $H\sim K$ 的非退化哈密顿量。回忆以下基本事实：
\begin{enumerate}[label=(\arabic*)]
 \item 哈密顿函数 $H$ 的非退化性只依赖其时间一映射 $\phi=\phi_H^1$；
 \item $\operatorname{Spec}(H)\subset\mathbb R$，即作用量泛函 $\mathcal A_H$ 的
 临界值集，是零测集（见 (Oh02) 的引理 2.2）；
 \item 若 $H,H'\mapsto\phi$ 且 $H\sim H'$，并且二者已归一化，则
 $\operatorname{Spec}(H)=\operatorname{Spec}(H')$（命题 18.3.6）；
 \item 函数 $H\mapsto\rho(H;a)$ 关于 $C_0^\infty(S^1\times M)$ 上的光滑拓扑连续
 （命题~\ref{prop:21-spectral-continuity}）；
 \item 从连通空间（例如区间 $[0,1]$）到 $\mathbb R$ 且值落在某个零测子集中的连续
 函数，只能是常函数。
\end{enumerate}

因 $H\sim K$，存在光滑族 $\mathcal H=\{H(s)\}_{0\leq s\leq1}$，使
$H(0)=H,H(1)=K$。定义 $\rho:[0,1]\to\mathbb R$ 为
$\rho(s)=\rho(H(s);a)$。由于所有 $s$ 的时间一映射均为
$\phi_{H(s)}^1=\phi_H^1$，$H(s)$ 都非退化；其像落在与 $s$ 无关的\emph{固定}子集
$\operatorname{Spec}(h)\subset\mathbb R$ 中，其中 $h=[H]=[K]$ 是道路同伦类。
由上述第 (2) 项，该集为零测集，因而全不连通。函数 $\rho$ 又由 $C^0$-连续性公理
连续，故必为常函数，于是
$\rho(H;a)=\rho(0)=\rho(1)=\rho(K;a)$，证明了非退化情形。
\footnote{译者注：原文此处回指第 (1) 项，但“作用谱是零测集”是上列第 (2) 项；
译文改正编号。}
\footnote{译者注：这里 $\rho$ 已被定义为参数函数 $s\mapsto\rho(H(s);a)$；原文等式
同时写 $\rho(H;a)$ 与 $\rho(0)$，记号层次混用，译文照录。}

现设 $H$ 退化。要强调的是，目前尚不知道退化哈密顿量是否满足谱性公理，因此不能
把上述非退化证明方案直接用于退化情形。改用基于下述引理的逼近论证。

\begin{lemma}
\label{lem:21-common-morse-perturbation}
设未必非退化的 $H\sim K$。则存在 $C^2$-收敛到 $0$ 的莫尔斯函数序列 $f_i$，使
$H\#f_i\sim K\#f_i$，且 $H\#f_i,K\#f_i$ 均非退化。
\end{lemma}

\begin{proof}
留作习题。
\end{proof}

应用此引理，
\[
 \rho(H;a)=\lim_{i\to\infty}\rho(H\#f_i;a)
 =\lim_{i\to\infty}\rho(K\#f_i;a)=\rho(K;a).
\]
第一种证明完成。

\stfhIndex{I147}{Zuiyouyantuo}{最优延拓（optimal continuation）}
\subsection{使用最优延拓的证明}

回忆同伦不变性的设定。给定 $H_0,H_1$，满足
$\phi_{H_0}^1=\phi_{H_1}^1$ 且 $[H_0]=[H_1]$。按定义，存在二参数族
$\{\phi(s,t)\}\subset\operatorname{Ham}(M,\omega)$，$(s,t)\in[0,1]^2$。存在
$\psi\in\operatorname{Ham}(M,\omega)$，使
\begin{equation}
 \begin{aligned}
 \phi(s,0)&\equiv\mathrm{id},&\quad \phi(s,1)&\equiv\psi
 &&\text{对所有 }s\in[0,1],\\
 \phi(0,t)&=\phi_{H_0}^t,& \phi(1,t)&=\phi_{H_1}^t
 &&\text{对所有 }t\in[0,1].
 \end{aligned}
 \label{eq:21-two-parameter-isotopy-boundary}
\end{equation}
记 $\mathcal H=H(s,t,x)$，其中 $H^s=H(s,\,\cdot\,,\,\cdot\,)$ 是生成向量场
\[
 \frac{\partial\phi}{\partial t}\circ\phi^{-1}
\]
的哈密顿量；又以 $F=F(s,t,x)$ 表示生成相应 $s$-哈密顿向量场
\[
 \frac{\partial\phi}{\partial s}\circ\phi^{-1}
\]
的哈密顿量。由~\eqref{eq:21-two-parameter-isotopy-boundary}，在 $t=0,1$ 附近沿
$t$ 方向、在 $s=0,1$ 附近沿 $s$ 方向重参数化
$\phi:[0,1]^2\to\operatorname{Ham}(M,\omega)$ 后，可设
\[
 F\equiv0\quad\text{于}\quad
 ([0,\epsilon]\times[0,1])\cup([1-\epsilon,1]\times[0,1]),
\]
特别地，$s$-哈密顿向量场 $X_F$ 对每个 $s$ 都关于 $t$ 一周期。

对给定 $[z^-,w^-]$、$[z^+,w^+]$，且
$[w^-\#u]=[w^+]\in\pi_2(z^+)$，考虑扰动柯西--黎曼方程
\begin{equation}
 \begin{cases}
 \partial u/\partial\tau-X_{F^\rho}(u)
 +J(\partial u/\partial t-X_{H^\rho}(u))=0,\\
 u(-\infty)=z^-,\quad u(\infty)=z^+,
 \end{cases}
 \label{eq:21-optimal-continuation-equation}
\end{equation}
并寻找有限能量解，其中
\[
 E_{(\phi^\rho,J)}(u):=\frac12\int_{\mathbb R\times[0,1]}
 \left(
 \left\lVert\frac{\partial u}{\partial\tau}-X_{F^\rho}(u)\right\rVert_J^2
 +\left\lVert\frac{\partial u}{\partial t}-X_{H^\rho}(u)\right\rVert_J^2
 \right).
\]
此方程等价于 $(du+P_{(K,J)}(u))^{0,1}=0$，其中
$K=H^\rho dt+F^\rho d\tau$。由于 $(H,F)$ 来自二参数族 $\phi$，有 $R_K\equiv0$。
又因 $F$ 在 $s=0,1$ 附近为 $0$，对充分大的 $R>0$，当 $|\tau|>R$ 时
$F^\rho\equiv0$。特别地，$\tau<-R$ 时~\eqref{eq:21-optimal-continuation-equation}
成为
\[
 \frac{\partial u}{\partial\tau}
 +J\left(\frac{\partial u}{\partial t}-X_{H_0}(u)\right)=0,
\]
而 $\tau>R$ 时成为
\[
 \frac{\partial u}{\partial\tau}
 +J\left(\frac{\partial u}{\partial t}-X_{H_1}(u)\right)=0.
\]

考察零维模空间以定义映射
\[
 h_{(\phi^\rho,J)}:CF(H_0,J_0)\longrightarrow CF(H_1,J_1),
\]
再如前考察一维模空间以证明它是链映射。还要分析该链映射如何影响弗洛尔复形的
滤过；为此，\emph{当}~\eqref{eq:21-optimal-continuation-equation} 有\emph{有限}能量解时，需研究
作用量差
\[
 \mathcal A_{H_0}([z^-,w^-])-
 \mathcal A_{H_1}([z^+,w^-\#u]).
\]
先回忆基本恒等式
\begin{equation}
 R_K=\frac{\partial F^\rho}{\partial t}
 -\frac{\partial H^\rho}{\partial\tau}
 +\{F^\rho,H^\rho\}=0.
 \label{eq:21-flat-curvature-identity}
\end{equation}
这是因为对子 $(H^\rho,F^\rho)$ 来自二参数同痕
$\phi=\phi(s,t)\in\operatorname{Ham}(M,\omega)$。对
\eqref{eq:21-optimal-continuation-equation} 应用命题
\ref{prop:20-energy-curvature-identity}，并用 $R_K=0$，得到
\[
 \mathcal A_{H_1}([z^+,w^-\#u])-\mathcal A_{H_0}([z^-,w^-])
 =-E_{(\phi^\rho,J)}(u)\leq0.
\]

\begin{proposition}
\label{prop:21-optimal-continuation-action}
设 $\phi=\{\phi(s,t)\}$ 及其相应对子 $(H,F)$ 如上。令
$\rho:\mathbb R\to[0,1]$ 光滑，且对某个 $R>0$，在 $\tau\leq-R$ 时
$\rho\equiv0$，在 $\tau\geq R$ 时 $\rho\equiv1$。给定
$[z^-,w^-]\in\operatorname{Crit}\mathcal A_{H_0}$，设
$u:\mathbb R\times[0,1]\to M$ 是~\eqref{eq:21-optimal-continuation-equation} 的解，
满足 $u(-\infty)=z^-,u(\infty)=z^+$。则
\begin{equation}
 \mathcal A_{H_1}([z^+,w^-\#u])
 \leq\mathcal A_{H_0}([z^-,w^-]).
 \label{eq:21-optimal-continuation-action-bound}
\end{equation}
\end{proposition}

立即得到：

\begin{corollary}
\label{cor:21-optimal-continuation-filtration}
链映射 $h_{(\phi;J)}:CF(H_0,J_0)\to CF(H_1,J_1)$ 保持滤过，即对所有
$\lambda\in\mathbb R$ 诱导链映射
\[
 h_{(\phi;J)}^\lambda:CF^\lambda(H_0,J_0)
 \longrightarrow CF^\lambda(H_1,J_1).
\]
\end{corollary}

显然，$s$-时间反演同痕
\[
 \widetilde\phi(s,t):=\phi(1-s,t)
\]
诱导保滤过映射
$h_{(\widetilde\phi;\widetilde J)}:
CF^\lambda(H_1,J_1)\to CF^\lambda(H_0,J_0)$。

现在完成 $\rho(H_0;a)=\rho(H_1;a)$ 的证明。只须证
$\rho(H_0;a)\leq\rho(H_1;a)$，反向同理，只需考虑 $s$-时间反演。给定
$\epsilon>0$，存在弗洛尔链 $\alpha_{H_1}$，使
$[\alpha_{H_1}]=a^\flat$ 且
$\lambda_{H_1}(\alpha_{H_1})\leq\rho(H_1;a)+\epsilon$。于是
\[
 \lambda_{H_0}\bigl(h_{(\widetilde\phi;\widetilde J)}(\alpha_{H_1})\bigr)
 \leq\lambda_{H_1}(\alpha_{H_1})\leq\rho(H_1;a)+\epsilon.
\]
又因 $h_{(\widetilde\phi;\widetilde J)}(\alpha_{H_1})\in CF_*(H_0)$ 且其同调类为
$a^\flat$，
\[
 \rho(H_0;a)
 \leq\lambda_{H_0}\bigl(h_{(\widetilde\phi;\widetilde J)}(\alpha_{H_1})\bigr)
 \leq\rho(H_1;a)+\epsilon.
\]
令 $\epsilon\to0$ 即得 $\rho(H_0;a)\leq\rho(H_1;a)$。这证明了哈密顿微分同胚
$\psi=\phi_{H_0}^1=\phi_{H_1}^1$ 非退化时的结论。一般情形由
$H\mapsto\rho(H;a)$ 的 $C^0$-连续性推出。
